% meta.concepts: fluid statics
% meta.tags: realistic
% acknowledge: Engineering Statics: Open and Interactive
% source: Chapter 9 Exercises of online edition (access: April 10, 2023)

\noindent A horizontal force $P$ is applied to a $W$ lbs box resting on an $32^\circ$ incline.  The line of action of $P$ passes through the center of gravity of the box.  The coefficients of friction between the box and the surface are $\mu_s = 0.1$ and $\mu_k = 0.080$.  The magnitude of force $P$ is such that the box is about the begin sliding up the incline.

\noindent At this condition, draw the free-body diagram (FBD) for the box in two ways:
\begin{enumerate}
  \item Showing the normal ($N$) and friction ($F$) forces as separate vectors.
  \item Combine the normal and friction forces into a single resultant vector ($R$) acting at a known angle.
\end{enumerate}

\noindent There is no need to compute the magnitude of the forces.  However, make sure all force vectors in either FBD have clear directions with numeric angles.

\begin{figure}[ht!]
  \centering
  \includegraphics[width=0.3\textwidth,height=0.5\textheight,keepaspectratio]{fig.png}
\end{figure}

\textit{Useful Relationships:}
\begin{itemize}
  \item Static Friction: $F_m = N \mu_s$
  \item Kintetic friction: $F_k = N \mu_k$
  \item Angle of static friction: $\tan \phi_s = \mu_s$
  \item Angle of kinetic friction: $\tan \phi_k = \mu_k$
\end{itemize}
